% Original source: Zhou, physical PDF pp. 178--182 (printed pp. 162--166).
% Figures: none.

% ZHOU-SOURCE-PAGE: 178 PRINTED: 162

$S_1$ and $S_2$ are geometrical Brownian motions, $f=\dfrac{S_2}{S_1}$ is a
geometric Brownian motion as well. One intuitive explanation is that both $\ln S_1$
and $\ln S_2$ follow normal distributions, so $\ln f=\ln S_2-\ln S_1$ follows a
normal distribution as well and $f$ follows a lognormal distribution. More rigorously,
we can apply Ito's lemma to $f=\dfrac{S_2}{S_1}$:

\[
\frac{\partial f}{\partial S_1}=-\frac{S_2}{S_1^2},\quad
\frac{\partial f}{\partial S_2}=\frac{1}{S_1},\quad
\frac{\partial^2 f}{\partial S_1^2}=\frac{2S_2}{S_1^3},\quad
\frac{\partial^2 f}{\partial S_2^2}=0,\quad
\frac{\partial^2 f}{\partial S_1\partial S_2}=-\frac{1}{S_1^2}.
\]

\[
\begin{aligned}
df={}&\frac{\partial f}{\partial S_1}dS_1+\frac{\partial f}{\partial S_2}dS_2
+\frac{1}{2}\frac{\partial^2 f}{\partial S_1^2}(dS_1)^2
+\frac{1}{2}\frac{\partial^2 f}{\partial S_2^2}(dS_2)^2
+\frac{\partial^2 f}{\partial S_1\partial S_2}dS_1dS_2\\
={}&-\mu_1\frac{S_2}{S_1}dt-\sigma_1\frac{S_2}{S_1}dW_{t,1}
+\mu_2\frac{S_2}{S_1}dt+\sigma_2\frac{S_2}{S_1}dW_{t,2}
+\sigma_1^2\frac{S_2}{S_1}dt-\rho\sigma_1\sigma_2\frac{S_2}{S_1}dt\\
={}&\left(\mu_2-\mu_1+\sigma_1^2-\rho\sigma_1\sigma_2\right)fdt
-\sigma_1f\,dW_{t,1}+\sigma_2f\,dW_{t,2}\\
={}&\left(\mu_2-\mu_1+\sigma_1^2-\rho\sigma_1\sigma_2\right)fdt
+\sqrt{\sigma_1^2-2\rho\sigma_1\sigma_2+\sigma_2^2}\times f\,dW_{t,3}.
\end{aligned}
\]

To make $f=\dfrac{S_2}{S_1}$ a martingale, set
$\mu_2-\mu_1+\sigma_1^2-\rho\sigma_1\sigma_2=0$ and we have

\[
\tilde E\left[\frac{S_{T,2}}{S_{T,1}}\right]=\frac{S_2}{S_1},
\]

and $\dfrac{S_{T,2}}{S_{T,1}}$ is a martingale under the new measure. The value of the
exchange option using $S_1$ as the numeraire is

\[
C_s=\tilde E\left[\max\left(\frac{S_{T,2}}{S_{T,1}}-1,0\right)\right],
\]

which is just the value of a call option with underlying asset price $S=\dfrac{S_2}{S_1}$,
strike price $K=1$, interest rate $r=0$, and volatility

\[
\sigma_s=\sqrt{\sigma_1^2-2\rho\sigma_1\sigma_2+\sigma_2^2}.
\]

So its value is $C_s=\dfrac{S_2}{S_1}N(d_1)-N(d_2)$, where

\[
d_1=\frac{\ln(S_2/S_1)+0.5\sigma_s^2\tau}{\sigma_s\sqrt{\tau}}
\quad\text{and}\quad d_2=d_1-\sigma_s\sqrt{\tau}.
\]

The payoff of the exchange option expressed in local currency is
$S_1C_s=S_2N(d_1)-S_1N(d_2)$.

% ZHOU-SOURCE-PAGE: 179 PRINTED: 163

\subsection{Other Finance Questions}

Besides option pricing problems, a variety of other quantitative finance problems are
tested in quantitative interviews as well. Many of these problems tend to be
position-specific. For example, if you are applying for a risk management job, prepare to
answer questions about VaR; for fixed-income jobs, get ready to answer questions about
interest rate models. As I explained in Chapter 1, it always helps if you grasp the basic
knowledge before the interview. In this section, we use several examples to show some
typical interview problems.

\subsubsection*{Portfolio optimization}

You are constructing a simple portfolio using two stocks $A$ and $B$. Both have the same
expected return of 12\%. The standard deviation of $A$'s return is 20\% and the standard
deviation of $B$'s return is 30\%; the correlation of their returns is 50\%. How will you
allocate your investment between these two stocks to minimize the risk of your portfolio?

\solution
Portfolio optimization has always been a crucial topic for investment management firms.
Harry Markowitz's mean-variance portfolio theory is by far the most well-known and
well-studied portfolio optimization model. The essence of the mean-variance portfolio
theory assumes that investors prefer (1) higher expected returns for a given level of
standard deviation/variance and (2) lower standard deviations/variances for a given level
of expected return. Portfolios that provide the minimum standard deviation for a given
expected return are termed efficient portfolios. The expected return and the variance of a
portfolio with $N$ assets can be expressed as

\[
\mu_p=w_1\mu_1+w_2\mu_2+\cdots+w_N\mu_N=w^T\mu.
\]

\[
\operatorname{var}(r_p)=\sum_{i=1}^{N}\sigma_i^2w_i^2
+\sum_{i\ne j}\sigma_{ij}w_iw_j=w^T\Sigma w.
\]

where $w_i$, $\forall i=1,\cdots,N$, is the weight of the $i$-th asset in the portfolio;
$\mu_i$, $\forall i=1,\cdots,N$, is the expected return of the $i$-th asset; $\sigma_i^2$
is the variance of $i$-th asset's return; $\sigma_{ij}=\rho_{ij}\sigma_i\sigma_j$ is the
covariance of the returns of the $i$-th and the $j$-th assets and $\rho_{ij}$ is their
correlation; $w$ is an $N\times1$ column vector of $w_i$'s; $\mu$ is an $N\times1$
column vector of $\mu_i$'s; $\Sigma$ is the covariance matrix of the returns of $N$
assets, an $N\times N$ matrix.

Since the optimal portfolio minimizes the variance of the return for a given level of
expected return, the efficient portfolio can be formulated as the following optimization
problem:

% ZHOU-SOURCE-PAGE: 180 PRINTED: 164

\[
\min_{w}\quad w^T\Sigma w,
\]

where $e$ is an $N\times1$ vector with all elements equal to 1.\footnote[13]{The
optimal weights have closed form solution $w^*=\lambda\Sigma^{-1}e+\gamma\Sigma^{-1}\mu$,
where $\lambda=\dfrac{C-\mu_pB}{D}$, $\gamma=\dfrac{\mu_pA-B}{D}$,
$A=e^T\Sigma^{-1}e>0$, $B=e^T\Sigma^{-1}\mu$, $C=\mu^T\Sigma^{-1}\mu>0$,
$D=AC-B^2$.}

\[
\textit{s.t.}\quad w^T\mu=\mu_p,\quad w^Te=1.
\]

For this specific problem, the expected returns are 12\% for both stocks. So $\mu_p$ is
always 12\% no matter what $w_A$ and $w_B$ ($w_A+w_B=1$) are. The variance of the
portfolio is

\[
\begin{aligned}
\operatorname{var}(r_p)&=\sigma_A^2w_A^2+\sigma_B^2w_B^2
+2\rho_{A,B}\sigma_A\sigma_Bw_Aw_B\\
&=\sigma_A^2w_A^2+\sigma_B^2(1-w_A)^2
+2\rho_{A,B}\sigma_A\sigma_Bw_A(1-w_A).
\end{aligned}
\]

Taking the derivative of $\operatorname{var}(r_p)$ with respect to $w_A$ and setting it
to zero, we have

\[
\frac{\partial\operatorname{var}(r_p)}{\partial w_A}
=2\sigma_A^2w_A-2\sigma_B^2(1-w_A)
+2\rho_{A,B}\sigma_A\sigma_B(1-w_A)
-2\rho_{A,B}\sigma_A\sigma_Bw_A=0.
\]

\[
\Rightarrow\quad
w_A=\frac{\sigma_B^2-\rho_{A,B}\sigma_A\sigma_B}
{\sigma_A^2-2\rho_{A,B}\sigma_A\sigma_B+\sigma_B^2}
=\frac{0.09-0.5\times0.2\times0.3}
{0.04-2\times0.5\times0.2\times0.3+0.09}=\frac{6}{7}.
\]

So we should invest 6/7 of the money in stock $A$ and 1/7 in stock $B$.

\subsubsection*{Value at risk}

Briefly explain what VaR is. What is the potential drawback of using VaR to measure the
risk of derivatives?

\solution
Value at Risk (VaR) and stress test---or more general scenario analysis---are two
important aspects of risk management. In the \emph{Financial Risk Manager Handbook},
\footnote[14]{\emph{Financial Risk Manager Handbook} by Phillippe Jorion is a comprehensive
book covering different aspects of risk management. A classic book for VaR is \emph{Value
at Risk}, also by Phillippe Jorion.} VaR is defined as the following: VaR is the maximum
loss over a target horizon such that there is a low, pre-specified probability that the
actual loss will be larger.

Given a confidence level $\alpha\in(0,1)$, the VaR can be implicitly defined as

\[
\alpha=\int_{\mathrm{VaR}}^{\infty}f(x)\,dx,
\]

where $x$ is the dollar profit (loss) and $f(x)$ is its probability density function. In
practice, $\alpha$ is often set to 95\% or 99\%. VaR is an extremely popular choice in
financial risk management since it summarizes the risk to a single dollar number.

% ZHOU-SOURCE-PAGE: 181 PRINTED: 165

Mathematically, it is simply the (negative) first or fifth percentile of the profit
distribution.

As a percentile-based measure on the profit distribution, VaR does not depend on the
shape of the tails before (and after) probability $1-\alpha$, so it does not describe the
loss on the left tail. When the profit/loss distribution is far from a normal distribution,
as in the cases of many derivatives, the tail portion has a large impact on the risk, and
VaR often does not reflect the real risk.\footnote[15]{Stress test is often used as a
complement to VaR by estimating the tail risk.} For example, let's consider a short
position in a credit default swap. The underlying asset is bond $A$ with a \$1M notional
value. Further assume that $A$ has a 3\% default probability and the loss given default is
100\% (no recovery). Clearly we are facing the credit risk of bond $A$. Yet if we use 95\%
confidence level, $VaR(A)=0$ since the probability of default is less than 5\%.

Furthermore, VaR is not sub-additive and is not a coherent measure of risk, which means
that when we combine two positions $A$ and $B$ to form a portfolio $C$, we do not always
have $VaR(C)\leq VaR(A)+VaR(B)$. For example, if we add a short position in a credit
default swap on bond $B$ with a \$1M notional value. $B$ also has a 3\% default
probability independent of $A$ and the loss given default is 100\%. Again we have
$VaR(B)=0$. When $A$ and $B$ form a portfolio $C$, the probability that at least one bond
will default becomes $1-(1-3\%)(1-3\%)\approx5.9\%$. So $VaR(C)=\$1M>VaR(A)+VaR(B)$.
Lack of sub-additivity clearly contradicts the intuitive idea that diversification reduces
risk. So it is a theoretical drawback of VaR.

(Sub-additivity is one property of a coherent risk measure. A risk measure $\rho(X)$ is
considered coherent if the following conditions holds: $\rho(X+Y)\leq\rho(X)+\rho(Y)$;
$\rho(aX)=a\rho(X)$, $\forall a>0$; $\rho(X)\leq\rho(Y)$, if $X\leq Y$; and
$\rho(X+k)=\rho(X)-k$ for any constant $k$. It is defined in \emph{Coherent Measure of
Risk} by Artzner, P., et al., Mathematical Finance, 9 (3):203--228. Conditional VaR is
a coherent risk measure.)

\subsubsection*{Duration and convexity}

The duration of a bond is defined as $D=-\frac{1}{P}\frac{\partial P}{\partial y}$, where
$P$ is the price of the bond and $y$ is yield to maturity. The convexity of a bond is
defined as $C=\frac{1}{P}\frac{\partial^2P}{\partial y^2}$. Applying Taylor's expansion,
$\dfrac{\Delta P}{P}\approx-D\Delta y+\dfrac{1}{2}C\Delta y^2$. when $\Delta y$ is
small, $\dfrac{\Delta P}{P}\approx-D\Delta y$.

For a fixed-rate bond with coupon rate $c$ and time-to-maturity $T$:

% ZHOU-SOURCE-PAGE: 182 PRINTED: 166

\[
T\uparrow\Rightarrow D\uparrow\quad c\uparrow\Rightarrow D\downarrow\quad y\uparrow\Rightarrow D\downarrow\qquad
T\uparrow\Rightarrow C\uparrow\quad c\uparrow\Rightarrow C\downarrow\quad y\uparrow\Rightarrow C\downarrow.
\]

Another important concept is dollar duration: $\$D=-\dfrac{dP}{dy}=P\times D$. Many
market participants use a concept called DV01: $DV01=-\dfrac{dP}{10,000\times dy}$,
which measures the price change when the yield changes by one basis point. For some bond
derivatives, such as swaps, dollar duration is especially important. A swap may have value
$P=0$, in which case dollar duration is more meaningful than duration.

When $n$ bonds with values $P_i$, $i=1,\cdots,n$, and Durations $D_i$ (convexities $C_i$)
form a portfolio, the duration of the portfolio is the value-weighted average of the
durations of the components:

\[
D=\sum_{i=1}^{n}\frac{P_i}{P}D_i\quad
\left(C=\sum_{i=1}^{n}\frac{P_i}{P}C_i\right),
\]

where $P=\sum_{i=1}^{n}P_i$. The dollar duration of the portfolio is simply the sum of
the dollar durations of the components: $\$D=\sum_{i=1}^{n}\$D_i$.

What are the price and duration of an inverse floater with face value \$100 and annual
coupon rate 30\%--3r that matures in 5 years? Assume that the coupons are paid
semiannually and the current yield curve is flat at 7.5\%.

\solution
The key to solving basic fixed-income problems is cash flow replication. To price a
fixed-income security with exotic structures, if we can replicate its cash flow using a
portfolio of fundamental bond types such as fixed-rate coupon bonds (including zero-coupon
bonds) and floating-rate bonds, no-arbitrage arguments give us the following conclusions:

Price of the exotic security = Price of the replicating portfolio

Dollar duration of the exotic security = Dollar duration of the replicating portfolio

To replicate the described inverse floater, we can use a portfolio constructed by shorting
3 floating rate bonds, which is worth \$100 each, and longing 4 fixed-rate bonds with a
7.5\% annual coupon rate, which is worth \$100 each as well. The coupon rate of a
floating-rate bond is adjusted every 0.5 years payable in arrear: the coupon rate paid at
$t+0.5y$ is determined at $t$. The cash flows of both positions and the whole portfolio are
summarized in the following table. It is apparent that the total cash flows of the portfolio
are the same as the described inverse floater. So the price of the inverse float is the price
of the replicating portfolio: $P_{\mathrm{inverse}}=\$100$.
